% Part: intuitionistic-logic
% Chapter: tableaux
% Section: introduction

\documentclass[../../../include/open-logic-section]{subfiles}

\begin{document}

\olfileid{int}{tab}{int}

\olsection{Pambuka}

!!^{tableau}s iku wit-wit tartamtu (kang cabange mudhun) saka !!{signed
  formula}s, yaiku pasangan kang dumadi saka tandha nilai kabeneran
($\True$ utawa $\False$) lan !!a{sentence}
\[
\sFmla{\True}{!A} \text{ utawa } \sFmla{\False}{!A}.
\]
!!^a{tableau} diwiwiti kanthi sawetara \emph{asumsi}. Saben
!!{signed formula} sabanjure diasilake kanthi ngetrapake salah siji
aturan inferensi. Sawetara aturan inferensi nambahake siji utawa
luwih !!{signed formula}s marang pucuk wit; liyane nambahake loro
pucuk anyar, saengga dadi rong cabang. Aturan ngasilake
!!{signed formula}s kanthi !!{formula} kang luwih prasaja tinimbang
ing !!{signed formula} kang dadi papan aturan iku ditrapake.
Yen sawijining cabang ngemot $\sFmla{\True}{!A}$ lan
$\sFmla{\False}{!A}$ bebarengan, cabang iku diarani \emph{katutup}.
Yen saben cabang ing !!a{tableau} katutup, kabeh !!{tableau} iku
katutup. !!{tableau} katutup iku !!a{derivation} kang nuduhake yen
himpunan !!{signed formula}s kang digunakake kanggo miwiti
!!{tableau} ora bisa dipenuhi. Iki bisa digunakake kanggo nemtokake
relasi $\Proves$: $\Gamma \Proves !A$ yen lan mung yen ana himpunan
winates~$\Gamma_0 = \{!B_1, \dots, !B_n\} \subseteq \Gamma$ kanthi
ana !!{tableau} katutup kanggo asumsi
\[
\{\sFmla{\False}{!A}, \sFmla{\True}{!B_1}, \dots, \sFmla{\True}{!B_n}\}.
\]

Kanggo logika intuisionistik, kita kudu nggedhekake pangerten
!!{signed formula} lan nyetel aturan kanggo konektif. Saliyane
tandha ($\True$ utawa $\False$), !!{formula}s ing !!{tableau}s
intuisionistik uga nduweni \emph{prefiks}~$\sigma$. Prefiks iku
urutan kang ora kosong saka wilangan bulat positif, yaiku
$\sigma \in (\PosInt)^* \setminus \{\emptyseq\}$.
Kita nulis prefiks kaya mangkono tanpa $\tuple{\ }$ kang ngubengi,
lan misahake saben !!{element}s nganggo~$.$ tinimbang $,$.
Yen $\sigma$ iku prefiks, mula $\sigma.n$ yaiku
$\sigma \concat \tuple{n}$; upamane yen $\sigma = 1.2.1$,
mula $\sigma.3$ yaiku $1.2.1.3$. Dadi, upamane,
\[
\sFmla{\True}{!A \lif (!B \lif !C)}[1.2]
\]
iku \emph{!!{signed formula} mawa prefiks} (utawa cukup
\emph{!!{formula} mawa prefiks} supaya ringkes).

Sacara intuitif, prefiks menehi jeneng donya ing model kang bisa
nyukupi !!{formula}s ing sawijining cabang !!a{tableau}, lan yen
$\sigma$ menehi jeneng sawijining donya, mula $\sigma.n$ menehi
jeneng donya kang bisa diakses saka (donya kang dijenengi dening)
~$\sigma$.

Ing model intuisionistik, relasi aksesibilitas refleksif lan transitif.
Tumrap prefiks, tegese $\sigma$ bisa diakses saka $\sigma$ dhewe,
lan uga saben prefiks kang nambahi~$\sigma$, yaiku saben prefiks
kang awujud $\sigma.n_1.\cdots.n_k$. Ayo nganggo notasi $\sigma.*$
kanggo nuduhake $\sigma$ dhewe lan saben ekstensine. Kanthi tembung
liya, prefiks $\sigma.*$ iku kabeh lan mung prefiks kang bisa
diakses saka~$\sigma$.

\end{document}
